\chapter{Certificate Validity and Blackout Safe-Hold}
\label{ch:cert-half-life-blackout}

This chapter re-states the battery blackout story in the form the repo can
actually defend.  The locked battery result is a \emph{safe-hold} result,
not an empirically identified physical half-life law.  When communication
disappears after a certificate has been issued, the current battery stack
collapses to a zero-power hold and preserves true-state safety throughout the
measured blackout window.

\section{Why Blackout Stress Matters}

Battery controllers do not always lose safety because of a bad optimisation
decision.  They can also lose safety because communication disappears after a
safe action has already been issued.  In a blackout, the operational question
is not ``what action is optimal now?'' but rather ``what certificate object
remains valid, for how long, and what happens when validity collapses to a
conservative hold?''

\section{Current Experimental Protocol}

The publication-facing artifacts are
\texttt{reports/publication/certificate\_half\_life\_blackout.csv} and
\texttt{reports/publication/blackout\_half\_life.csv}.  The underlying
battery-domain protocol is:

\begin{enumerate}
  \item Generate a \texttt{drift\_combo} episode with the DC3S loop active.
  \item Advance the controller to a freeze point where a safe action has been
    certified.
  \item Hold the certified action fixed through blackout horizons of
    4, 12, and 24 hours in the publication-facing package.
  \item Evaluate the true-state battery trajectory for violations during the
    blackout.
\end{enumerate}

\section{Locked Results}

Table~\ref{tab:blackout-study} reproduces the battery-safe-hold surface from
the canonical publication artifact.

\begin{table}[ht]
\centering
\caption{Battery blackout safe-hold results from
  \texttt{reports/publication/certificate\_half\_life\_blackout.csv}.}
\label{tab:blackout-study}
\begin{tabular}{lrrrrr}
\toprule
\textbf{Blackout (h)} & \textbf{Frozen charge} & \textbf{Frozen discharge} &
  \textbf{TSVR (\%)} & \textbf{Sev.\ P95 (MWh)} & \textbf{Coverage} \\
\midrule
  0  & 0.0 & 0.0 & 0.0 & 0.00 &   0\% \\
  4  & 0.0 & 0.0 & 0.0 & 0.00 & 100\% \\
 12  & 0.0 & 0.0 & 0.0 & 0.00 & 100\% \\
 24  & 0.0 & 0.0 & 0.0 & 0.00 & 100\% \\
\bottomrule
\end{tabular}
\end{table}

The battery-operational interpretation is straightforward: the frozen
certificate currently resolves to a benign zero-power hold.  In the locked
publication surface, no measured blackout horizon produces a true-state
violation.

\section{Width-Expansion Proxy}

The companion artifact \texttt{reports/publication/blackout\_half\_life.csv}
tracks the \emph{certificate-width reaction} at blackout onset:

\begin{table}[ht]
\centering
\caption{Width-expansion proxy at blackout onset from
  \texttt{reports/publication/blackout\_half\_life.csv}.}
\label{tab:blackout-width-proxy}
\begin{tabular}{rrrr}
\toprule
\textbf{Baseline width} & \textbf{Peak width} & \textbf{Half-target width} &
  \textbf{Proxy half-life (steps)} \\
\midrule
376.87 & 5353.04 & 2864.95 & 0 \\
\bottomrule
\end{tabular}
\end{table}

This is a \emph{certificate-width proxy}, not an empirical physical-failure
surface.  The zero-step proxy value means the width reaction is immediate at
fault onset; it does \emph{not} mean the battery becomes unsafe immediately.

\section{Operational Consequence}

For the battery domain, this is still an important control result.  When
telemetry disappears entirely, the present DC3S stack has an auditable path
to a benign hold action rather than blindly persisting with a stale,
high-impact dispatch.  In asset terms, that means the blackout policy favors
state preservation over speculative useful work once the measurement channel
goes dark.

\section{What the Current Result Does and Does Not Show}

This locked result establishes a bounded claim:

\begin{itemize}
  \item The present battery implementation can enter a conservative frozen
    action that remains safe throughout the measured blackout window.
  \item The companion width proxy confirms that the certificate reacts
    immediately to blackout onset by inflating uncertainty.
  \item The current package does \emph{not} yet identify a finite empirical
    safety-breakdown horizon for a stale aggressive action, so the chapter
    should not be read as a final certificate-expiration law.
\end{itemize}

\section{Summary}

The blackout chapter now matches the locked evidence.  The defended battery
claim is that certificate validity currently degrades into a safe-hold path
with zero measured violations, while the half-life artifact should be read as
an interval-width proxy rather than as an empirical physical-failure clock.
